\documentclass[11pt,a4paper]{article}

\usepackage[utf8]{inputenc}
\usepackage[T1]{fontenc}
\usepackage{lmodern}
\usepackage{amsmath,amssymb,amsfonts}
\usepackage{booktabs}
\usepackage{tabularx}
\usepackage{geometry}
\usepackage{xcolor}
\usepackage{tcolorbox}
\usepackage{hyperref}

\geometry{margin=2.2cm}
\hypersetup{
    colorlinks=true,
    linkcolor=blue!70!black,
    citecolor=blue!70!black,
    urlcolor=blue!70!black
}

\title{\textbf{AEON-$T_{\mu\nu}$ \& Philosophy of Time: Neuroscience, Temporal Binding, and the Operational Simultaneity Horizon}}
\author{Luís Pimenta Damásio \\ \small Independent Researcher, Lisbon, Portugal \quad | \quad ORCID: 0009-0001-7606-7557}
\date{November 18, 2026}

\begin{document}

\maketitle

\begin{tcolorbox}[colback=gray!6, colframe=gray!65!black, title=\textbf{Epistemic Status \& AI Collaboration Disclosure}]
\small
\textbf{Epistemic Status:} This paper presents a speculative theoretical and philosophical synthesis extending the \textbf{AEON-$T_{\mu\nu}$} framework into systems neuroscience. Standard neurophysiological mechanisms (thalamocortical oscillations, axonal conduction delays, Perturbational Complexity Index) and open-system physics (GKSL/Lindblad master equations, Caldeira--Leggett thermal dissipation) are strictly distinguished from the proposed ontological overlays ($t_{\mathrm{pre}}$, $\rho_\tau$, $A_s^v$, and $L_{\mathrm{sim}}$), which remain unverified theoretical hypotheses subject to empirical falsification.\\[4pt]
\textbf{AI Collaboration Note:} AI assistance was utilized for mathematical typesetting, structural organization, cross-reference auditing, and document preparation.
\end{tcolorbox}

\begin{abstract}
\noindent The neurobiological binding problem asks how billions of spatially separated neurons, processing distinct sensory modalities across distributed cortical regions with finite axonal conduction delays, generate a single, temporally unified conscious present. Standard computational and neurobiological models explain local feature encoding and structural connectivity, yet often leave the temporal unity of lived experience either unexplained or treated as an emergent epiphenomenon. Building upon the \emph{Master Compendium}, the \emph{AEON-T14.html Computational Observatory}, and the \emph{Executive Manifesto and Front-Door Reader's Guide}, this paper specializes the \emph{Covariant Objective-Time Master Equation} to terrestrial biological scales. We model the awake thalamocortical network as a metabolically driven coherent whole functioning as an intrinsic clock ($t_{\mathrm{pre}}$) with continuous diachronic identity ($A_s^v$). By introducing an effective cortical operational simultaneity horizon ($L_{\mathrm{sim}}^{\mathrm{cortex}}$) governed by the competition between active oscillatory phase-locking and effective inter-areal phase slip ($\Gamma_{\mathrm{eff}}$), we show that unified temporal binding occurs when $L_{\mathrm{sim}}^{\mathrm{cortex}} > L_{\mathrm{cortex}}$ ($D_{\mathrm{loc}} \to 1$). Conversely, general anesthesia, deep slow-wave sleep, and pathological fragmentation correspond to a dynamical contraction of $L_{\mathrm{sim}}^{\mathrm{cortex}}$ below the cortical diameter, dissolving global phase-locking into decoupled local modules without requiring macroscopic quantum coherence in warm wet tissue. Finally, we formalize the epigenetics of consciousness—distinguishing the physical preconditions of temporal binding from the irreducible phenomenal plane of qualitative experience—and outline falsifiable empirical predictions.
\end{abstract}

\section{Introduction: The Temporal Binding Problem and the Cartesian Split}
A central challenge in systems neuroscience and the philosophy of mind is the \emph{binding problem}: how distributed neural assemblies—separated by centimeters of cortical tissue and communicating via finite axonal conduction velocities ($v_{\mathrm{ax}} \sim 1$--$100\;\mathrm{m/s}$)—combine distinct sensory attributes (color, motion, timbre, interoception) into a single, seamless experiential ``now'' \cite{treisman1996,singer1999,varela2001}. While anatomical tract-tracing and functional neuroimaging have mapped structural connectomes in remarkable detail, a deeper conceptual gap persists between external third-person descriptions of neural spikes occurring along a laboratory clock ($t_{\mathrm{obs}}$) and the first-person unity of lived temporal duration.

Mainstream approaches typically fall into two camps. Reductionist models treat the brain strictly as an \emph{object-as-object}: a passive network of biophysical switches whose temporal unity is assumed to be an epiphenomenal illusion generated by data integration. Conversely, non-reductive panpsychist or quantum-brain proposals (such as Orchestrated Objective Reduction \cite{hameroff2014}) attempt to rescue subjectivity either by attributing primitive experience to every elementary particle or by positing fragile, long-range quantum superpositions inside warm, noisy cellular environments \cite{tegmark2000}.

As established in the \emph{Executive Manifesto and Front-Door Reader's Guide}, the \textbf{AEON-$T_{\mu\nu}$} framework resolves this impasse through an ontological inversion grounded in two foundational principles:
\begin{enumerate}
    \item \textbf{Existence as Duration and Spatial Localization:} No physical entity exists instantaneously; to persist across a finite interval is to actively integrate change across a bounded spatial domain.
    \item \textbf{The Second Postulate of Time (Object-as-Subject):} Every dynamically coherent whole functions as an intrinsic internal clock ($t_{\mathrm{pre}}$) ticking at its own objective temporal intensity ($\rho_\tau$) and maintaining diachronic identity ($A_s^v$) so long as its internal phase coherence survives environmental scattering.
\end{enumerate}
In this paper, we apply this architecture to the human nervous system. Rather than invoking macroscopic quantum superpositions across the skull, we demonstrate how classical electromagnetic and metabolic phase-locking across thalamocortical loops establishes a macroscopic \emph{operational simultaneity horizon} ($L_{\mathrm{sim}}^{\mathrm{cortex}}$) that exceeds the physical diameter of the cortex ($L_{\mathrm{cortex}} \approx 0.15$--$0.20\;\mathrm{m}$), binding distributed neural ensembles into a single unified \emph{object-as-subject}.

\section{Specializing the Master Equation to Cortical Scales}
In the \emph{Executive Manifesto and Front-Door Reader's Guide}, open-system evolution across curved spacetime is governed by the \emph{Covariant Objective-Time Master Equation}:
\begin{equation}
    \frac{d\hat{\rho}}{dt_{\mathrm{pre}}} = -\frac{i}{\hbar}\left[\hat{H}_{\mathrm{eff}}\left(\langle \hat{T}_{\mu\nu}\rangle_{\mathrm{ren}}\right), \hat{\rho}\right] + \mathcal{D}_{\mathrm{th}}[\hat{\rho}] + \mathcal{D}_{\mathrm{grav}}[\hat{\rho}],
    \label{eq:master}
\end{equation}
where $dt_{\mathrm{pre}} = \sqrt{N^2(\theta)}\,dt$ denotes the intrinsic proper time increment of the coherent system, $\mathcal{D}_{\mathrm{th}}$ represents Caldeira--Leggett thermal dissipation \cite{caldeira1983}, and $\mathcal{D}_{\mathrm{grav}}$ represents gravitational self-energy decoherence \cite{penrose1996,diosi1989}. Across a spatial baseline $\Delta\ell$, off-diagonal phase coherence between separated nodes decays according to the spatial locality factor:
\begin{equation}
    \mathcal{D}_{\mathrm{loc}}(\Delta\ell) = \exp\!\left(-\frac{\Delta\ell}{L_{\mathrm{sim}}(r)}\right), \qquad \mathcal{S}_{\mathrm{sim}} = \Phi_\tau(\Omega)\,\mathcal{D}_{\mathrm{loc}}(\Delta\ell),
    \label{eq:dloc}
\end{equation}
where $\Phi_\tau(\Omega) = \Phi_{\mathrm{floor}} + (1 - \Phi_{\mathrm{floor}})e^{-\Omega}$ tracks timeline integrity against accumulated decoherence action $\Omega = \int \Gamma_{\mathrm{tot}}\,dt_{\mathrm{pre}}$, and $\mathcal{S}_{\mathrm{sim}}$ measures closed-loop operational synchronization fidelity.

\subsection{The Terrestrial Weak-Field Limit and Microscopic Decoherence}
At the terrestrial surface ($r = R_\oplus$) and biological scale ($L_{\mathrm{cortex}} \sim 0.18\;\mathrm{m}$, corresponding to conformal floor $\zeta = \log_{10}(L/\ell_P) \approx 34.05$ in the \emph{AEON-T14.html Computational Observatory}), spacetime curvature is weak ($N^2(\theta) \approx 1 - 2GM_\oplus/R_\oplus c^2 \approx 1 - 1.39 \times 10^{-9}$) and the semiclassical breakdown index sits far below unity ($\Xi_{\mathrm{QG}} \sim 10^{-9} \ll 1$). At physiological temperature ($T_{\mathrm{eff}} \approx 310.15\;\mathrm{K}$), microscopic quantum superpositions of individual ions or proteins undergo rapid thermal decoherence:
\begin{equation}
    \Gamma_{\mathrm{th}}^{\mathrm{micro}} \sim \frac{L^2 (k_B T_{\mathrm{eff}})^3}{\hbar^3 c^2} \gg 10^{20}\;\mathrm{s}^{-1}, \qquad \Phi_\tau^{\mathrm{micro}} \to \Phi_{\mathrm{floor}}.
\end{equation}
Consequently, as Tegmark \cite{tegmark2000} rightly emphasized, bare microscopic quantum wavefunctions cannot span the brain unscattered. However, in \textbf{AEON-$T_{\mu\nu}$}, macroscopic temporal unity does not require fragile microscopic quantum superpositions; it operates at the \emph{mesoscopic and macroscopic dynamical level} of collective neural oscillations.

\subsection{Effective Cortical Operational Simultaneity Horizon ($L_{\mathrm{sim}}^{\mathrm{cortex}}$)}
When metabolic free energy (ATP hydrolysis) drives populations of cortical pyramidal neurons into limit-cycle oscillations across nested frequency bands (theta $4$--$8\;\mathrm{Hz}$, alpha $8$--$12\;\mathrm{Hz}$, beta $13$--$30\;\mathrm{Hz}$, and gamma $30$--$100\;\mathrm{Hz}$), the relevant dynamical variables are the collective phase fields $\theta_i(t_{\mathrm{pre}})$ of local neural columns communicating at effective signal velocity $v_{\mathrm{sig}}$ (spanning fast ephaptic/electromagnetic field coupling and myelinated thalamocortical tracts).

Replacing bare microscopic thermal scattering with the \emph{effective inter-areal phase-slip rate} $\Gamma_{\mathrm{eff}}$—defined as the rate at which synaptic noise, metabolic fluctuations, and conduction dispersion disrupt phase-locking across cortical assemblies—Equation~(\ref{eq:dloc}) specializes to the \textbf{Cortical Operational Simultaneity Horizon}:
\begin{equation}
    L_{\mathrm{sim}}^{\mathrm{cortex}} = \frac{v_{\mathrm{sig}}}{2\,\Gamma_{\mathrm{eff}}}, \qquad \mathcal{D}_{\mathrm{loc}}^{\mathrm{cortex}}(\Delta\ell) = \exp\!\left(-\frac{\Delta\ell}{L_{\mathrm{sim}}^{\mathrm{cortex}}}\right).
    \label{eq:lsim_cortex}
\end{equation}
Equation~(\ref{eq:lsim_cortex}) yields a sharp, quantitative criterion for unified temporal binding:
\begin{itemize}
    \item \textbf{Bound Regime ($L_{\mathrm{sim}}^{\mathrm{cortex}} > L_{\mathrm{cortex}}$):} When active thalamocortical resonance suppresses $\Gamma_{\mathrm{eff}}$ below $v_{\mathrm{sig}} / (2 L_{\mathrm{cortex}})$, round-trip signaling lag $\Delta\tau_{\mathrm{sig}} = 2\Delta\ell / v_{\mathrm{sig}}$ remains shorter than the collective phase-coherence lifetime $\tau_{\mathrm{coh}} = 1/\Gamma_{\mathrm{eff}}$. Spatial locality decay approaches unity ($\mathcal{D}_{\mathrm{loc}}^{\mathrm{cortex}} \to 1$), locking distributed cortical areas into a single shared proper duration ($t_{\mathrm{pre}}$) and unbroken diachronic identity ($A_s^v$).
    \item \textbf{Unbound / Fragmented Regime ($L_{\mathrm{sim}}^{\mathrm{cortex}} \ll L_{\mathrm{cortex}}$):} When $\Gamma_{\mathrm{eff}}$ rises or long-range signaling $v_{\mathrm{sig}}$ is attenuated, $\mathcal{D}_{\mathrm{loc}}^{\mathrm{cortex}} \to 0$ across inter-areal distances. Distributed cortical modules continue to process local inputs independently, but global phase-locking collapses; the brain ceases to operate as a single unified \emph{object-as-subject} and reverts to an ensemble of decoupled local clocks.
\end{itemize}

\section{The Conformal Scale Bifurcation and Clinical States of Consciousness}
As detailed in the \emph{Executive Manifesto and Front-Door Reader's Guide}, riding the \emph{Conformal Preset Elevator} to terrestrial macroscopic scales ($\zeta \approx 31$--$35$) exposes a fundamental bifurcation: an inert grain of sand and a living mammalian brain occupy nearly the same conformal scale, yet exhibit opposite internal temporal architectures.

\begin{table}[htbp]
\centering
\caption{Dynamical Regimes of Cortical Temporal Binding Across States of Consciousness}
\label{tab:states}
\small
\begin{tabularx}{\textwidth}{@{}p{2.9cm}p{2.3cm}p{2.4cm}p{2.4cm}X@{}}
\toprule
\textbf{System / Brain State} & \textbf{Effective Phase Slip ($\Gamma_{\mathrm{eff}}$)} & \textbf{Horizon Ratio ($L_{\mathrm{sim}}^{\mathrm{cortex}} / L_{\mathrm{cortex}}$)} & \textbf{Sync Fidelity ($\mathcal{S}_{\mathrm{sim}}$)} & \textbf{Ontological \& Empirical Signature} \\
\midrule
\textbf{Inert Macro-Object} (e.g., Sand Grain) & $\Gamma_{\mathrm{eff}} \to \infty$ (passive thermal bath) & $\ll 10^{-10}$ & $\to 0$ & Pure \emph{object-as-object}; zero global phase-locking; statistical aggregate of micro-clocks. \\
\addlinespace
\textbf{Awake Conscious Perception} & Low (active ATP-driven resonance) & $> 1$ (global cortical span) & $\to 1$ (high) & Unified \emph{object-as-subject}; single $t_{\mathrm{pre}}$; high Perturbational Complexity Index ($\mathrm{PCI} > 0.31$) \cite{casali2013}. \\
\addlinespace
\textbf{REM Sleep / Dreaming} & Low–Moderate (internally driven loops) & $\gtrsim 1$ (internal thalamocortical) & Moderate–High (internal) & Internally bound \emph{object-as-subject} decoupled from external sensory channels ($\mathcal{D}_{\mathrm{loc}}^{\mathrm{ext}} \to 0$); high PCI. \\
\addlinespace
\textbf{General Anesthesia (Propofol / Xenon)} & High (GABAergic hyperpolarization) & $< 1$ (local columns only) & $\to 0$ (long-range) & Horizon contraction; local sensory responses persist, but global integration collapses ($\mathrm{PCI} < 0.31$) \cite{mashour2020}. \\
\addlinespace
\textbf{Generalized Tonic-Clonic Seizure} & Hypersynchronous lock (zero differentiation) & $\gg 1$ (pathological saturation) & Saturated / Rigid & Loss of informational differentiation; attentional condensation freezes; consciousness lost despite high synchrony. \\
\bottomrule
\end{tabularx}
\end{table}

\subsection{Anesthesia as Contraction of the Operational Simultaneity Horizon}
General anesthesia provides a direct empirical test of Equation~(\ref{eq:lsim_cortex}). Under propofol or sevoflurane, primary sensory cortices still respond to external stimuli—meaning local neural clocks continue to tick—yet conscious awareness vanishes. Neurophysiological recordings demonstrate that anesthetics selectively disrupt long-range frontoparietal feedback connectivity and phase synchronization \cite{mashour2020,lee2013}. In \textbf{AEON-$T_{\mu\nu}$}, anesthetic agents increase effective inter-areal phase slip $\Gamma_{\mathrm{eff}}$ and attenuate effective long-range signaling bandwidth, shrinking $L_{\mathrm{sim}}^{\mathrm{cortex}}$ below $L_{\mathrm{cortex}}$. The macroscopic \emph{object-as-subject} dissolves into isolated cortical islands, directly mirroring the empirical drop in the Perturbational Complexity Index (PCI) below the consciousness threshold ($\mathrm{PCI}^* \approx 0.31$) established by Casali et al.\ \cite{casali2013}.

\subsection{Why Generalized Seizures Abolish Consciousness: Differentiation Within $L_{\mathrm{sim}}$}
Table~\ref{tab:states} highlights a crucial subtlety: maintaining $L_{\mathrm{sim}}^{\mathrm{cortex}} > L_{\mathrm{cortex}}$ is a \emph{necessary} condition for unified temporal binding, but not a sufficient one. During a generalized epileptic seizure, hypersynchronous runaway firing locks the entire cortex into a single, rigid, low-dimensional oscillation. Although spatial phase-locking is maximal, internal informational differentiation collapses to zero: the system loses its capacity to condense distinct sensory or cognitive states from \emph{functional infinity}. Thus, conscious binding requires \emph{metastable phase-locking}—high spatial integration ($L_{\mathrm{sim}}^{\mathrm{cortex}} > L_{\mathrm{cortex}}$) paired with high dynamical repertoire differentiation \cite{tononi2016}.

\section{The Epigenetics of Consciousness and Attentional Condensation}
To ground neurobiology without falling into either naive materialism or dualistic mysticism, we incorporate the pre-mathematical metaphysics of consciousness articulated in the \emph{Executive Manifesto and Front-Door Reader's Guide}:
\begin{enumerate}
    \item \textbf{Consciousness as the Pre-Mathematical Plane of Duration:} At every waking moment, lived consciousness opens a continuous, pre-mathematical temporal plane—without sharp beginning, sharp end, or discrete clock-ticks in between.
    \item \textbf{Attentional Condensation from Functional Infinity:} Just as an uncoupled quantum wave inhabits \emph{functional infinity} ($\Phi_\tau = 1$) before condensing into a single actualized interaction ($\Phi_\tau \to \Phi_{\mathrm{floor}}$), the unbounded temporal plane of consciousness continually condenses into \emph{what is actively under focus}—never the whole infinity at once, but ``this'' or ``that'' intentional object populating the present.
    \item \textbf{Past and Future as Epistemic Projections:} When a conscious observer interacts with an external object and applies physical laws to extrapolate its trajectory backward or forward, those reconstructed pasts and anticipated futures are \emph{imagined projections} within consciousness, not pre-existing rooms inside a static four-dimensional block universe. Physically, only ongoing proper duration ($t_{\mathrm{pre}}$) and diachronic identity ($A_s^v$) exist.
    \item \textbf{Epigenetic Emergence and Dual Preconditions:} Consciousness behaves epigenetically with respect to matter: it is not a separate physical substance, yet it requires two strict physical preconditions to manifest: (i) an adequate biological substrate capable of sustaining an intrinsic clock (conformal floor $\zeta \approx 34$--$35$), and (ii) a special dynamical configuration ($L_{\mathrm{sim}}^{\mathrm{cortex}} > L_{\mathrm{cortex}}$ with metastable differentiation). When condition (ii) fails—as under general anesthesia—the conscious plane disappears even though the biological substrate remains intact.
\end{enumerate}

\section{Comparative Positioning Against Leading Theories of Consciousness}
Situating \textbf{AEON-$T_{\mu\nu}$} alongside established neuroscientific and philosophical frameworks clarifies its distinct explanatory contribution.

\begin{table}[htbp]
\centering
\caption{Comparative Matrix: Neural Binding and Consciousness Across Frameworks}
\label{tab:comp}
\small
\begin{tabularx}{\textwidth}{@{}p{2.7cm}p{2.8cm}p{2.9cm}X@{}}
\toprule
\textbf{Framework} & \textbf{Primary Mechanism} & \textbf{Treatment of Time} & \textbf{Contrast with AEON-$T_{\mu\nu}$} \\
\midrule
\textbf{Global Neuronal Workspace (GNW)} \cite{dehaene2011} & Non-linear frontoparietal ``ignition'' and broadcast. & Implicit external lab time ($t_{\mathrm{obs}}$). & GNW explains cognitive access and reportability; \textbf{AEON-$T_{\mu\nu}$} supplies the underlying physical metric ($L_{\mathrm{sim}}^{\mathrm{cortex}} > L_{\mathrm{cortex}}$) that unifies the workspace into a single $t_{\mathrm{pre}}$. \\
\addlinespace
\textbf{Integrated Information Theory (IIT)} \cite{tononi2016} & Irreducible cause--effect structure ($\Phi_{\mathrm{IIT}}$). & Discrete state transitions across causal steps. & IIT quantifies informational irreducibility; \textbf{AEON-$T_{\mu\nu}$} anchors integration in continuous covariant proper duration ($dt_{\mathrm{pre}}$) and signaling-bounded horizons. \\
\addlinespace
\textbf{Communication Through Coherence (CTC)} \cite{fries2015} & Rhythmic gamma/beta phase alignment gating synaptic gain. & Oscillatory phase relationships in lab time. & CTC provides the exact neurophysiological mechanism that lowers $\Gamma_{\mathrm{eff}}$ and expands $L_{\mathrm{sim}}^{\mathrm{cortex}}$ in our framework. \\
\addlinespace
\textbf{Orchestrated Objective Reduction (Orch-OR)} \cite{hameroff2014} & Gravitational collapse of tubulin quantum superpositions. & Discrete quantum collapse events ($\tau \approx \hbar / E_G$). & Unlike Orch-OR, \textbf{AEON-$T_{\mu\nu}$} does \emph{not} require fragile macroscopic quantum superpositions in warm microtubules; binding occurs via classical/metabolic phase-locking inside $L_{\mathrm{sim}}^{\mathrm{cortex}}$. \\
\addlinespace
\textbf{AEON-$T_{\mu\nu}$ (Object-as-Subject)} & Phase-locking inside $L_{\mathrm{sim}}^{\mathrm{cortex}} > L_{\mathrm{cortex}}$; attentional condensation. & Intrinsic proper duration ($t_{\mathrm{pre}}$) and diachronic identity ($A_s^v$). & Bridges subjective lived duration and objective biophysics while explicitly delimiting structural temporal binding from qualitative qualia. \\
\bottomrule
\end{tabularx}
\end{table}

\section{Explicit Structural Boundaries and Falsifiable Predictions}
Consistent with the methodological rigor established in the \emph{Executive Manifesto and Front-Door Reader's Guide}, we state four explicit boundaries where this framework stops:
\begin{enumerate}
    \item \textbf{The Hard Problem of Qualitative Qualia:} \textbf{AEON-$T_{\mu\nu}$} solves the \emph{structural and temporal binding problem}—explaining how distributed neural processes merge into a single unified temporal perspective ($t_{\mathrm{pre}}$) with continuous diachronic identity ($A_s^v$). It does \emph{not} claim to deduce why a specific oscillatory pattern feels qualitatively like the color red or the warmth of sunlight; the phenomenal plane itself remains epigenetic to its physical preconditions.
    \item \textbf{No Macroscopic Quantum Superposition Claim:} Because $\Gamma_{\mathrm{th}}^{\mathrm{micro}} \gg 10^{20}\;\mathrm{s}^{-1}$ at $310\;\mathrm{K}$, individual neural microtubules or ion channels sit at $\Phi_\tau^{\mathrm{micro}} \approx \Phi_{\mathrm{floor}}$. All macroscopic cortical coherence in this model is classical, electromagnetic, and metabolic.
    \item \textbf{Coarse-Graining of Cortical Boundaries ($L_{\mathrm{cortex}}$):} Defining the exact spatial boundary of the coherent whole requires specifying the active thalamocortical network engaged in a given task; subcortical autonomic loops operating outside $L_{\mathrm{sim}}^{\mathrm{cortex}}$ remain unconscious background clocks.
    \item \textbf{Falsifiable Empirical Predictions:}
    \begin{itemize}
        \item \emph{Prediction 1 (Horizon-to-Diameter Threshold):} Across graded anesthetic titration (e.g., propofol target-controlled infusion), the subjective loss of temporal unity and behavioral responsiveness must coincide with the quantitative crossing $L_{\mathrm{sim}}^{\mathrm{cortex}} = v_{\mathrm{sig}} / (2\Gamma_{\mathrm{eff}}) \approx L_{\mathrm{cortex}}$, measurable via combined high-density EEG/MEG phase-locking value decay and white-matter conduction tractography.
        \item \emph{Prediction 2 (Attentional Condensation Dynamics):} Shifts of conscious focus from broad diffuse awareness to narrow focal attention (``this'' or ``that'') will correspond to measurable cross-frequency phase-amplitude coupling (theta--gamma nesting) that transiently sharpens $\mathcal{D}_{\mathrm{loc}}^{\mathrm{cortex}} \to 1$ across task-relevant sensory--frontal subnetworks while decoupling unattended background assemblies.
    \end{itemize}
\end{enumerate}

\section{Conclusion and Roadmap Integration}
By specializing the \emph{Covariant Objective-Time Master Equation} to the human cortex, this paper demonstrates that the neurobiological binding problem is fundamentally a problem of \emph{operational simultaneity}. When active metabolic free energy suppresses effective inter-areal phase slip ($\Gamma_{\mathrm{eff}}$) such that the cortical operational simultaneity horizon exceeds the physical diameter of the brain ($L_{\mathrm{sim}}^{\mathrm{cortex}} > L_{\mathrm{cortex}}$), spatially distributed neural assemblies phase-lock into a single \emph{object-as-subject} ticking along one unified proper duration ($t_{\mathrm{pre}}$).

This neurobiological synthesis connects directly to the remaining investigations in our research program: extending diachronic identity ($A_s^v$) and metabolic entropy export across whole living organisms in our forthcoming study on \emph{biological persistence and metabolic coherence}; testing $\mathcal{D}_{\mathrm{loc}}(\Delta\ell)$ across precision fiber and satellite links in \emph{optical geodesy}; illustrating wave condensation in \emph{everyday illuminated spaces}; and completing the overarching philosophical synthesis in our treatise on \emph{the ontology of duration}.

\begin{thebibliography}{99}

\bibitem{damasio_comp1}
L.~Pimenta Damásio, \emph{Master Compendium — Foundations of Temporal Metrology and Operational Quantum Mechanics (v1.0)}, Zenodo (2026), DOI: \href{https://doi.org/10.5281/zenodo.24163322}{10.5281/zenodo.24163322}.

\bibitem{damasio_comp2}
L.~Pimenta Damásio, \emph{Master Compendium — Foundations of Temporal Metrology and Operational Quantum Mechanics (v2.0)}, Zenodo (2026), DOI: \href{https://doi.org/10.5281/zenodo.24190229}{10.5281/zenodo.24190229}.

\bibitem{damasio_obs}
L.~Pimenta Damásio, \emph{AEON-T14.html: Covariant Kerr--Newman QFTCS \& Semiclassical Back-Reaction Observatory}, Zenodo (2026), DOI: \href{https://doi.org/10.5281/zenodo.24244501}{10.5281/zenodo.24244501}.

\bibitem{damasio_manifesto}
L.~Pimenta Damásio, \emph{AEON-$T_{\mu\nu}$ \& Philosophy of Time: Executive Manifesto and Front-Door Reader's Guide}, Zenodo Preprint (2026).

\bibitem{treisman1996}
A.~Treisman, ``The binding problem,'' \emph{Current Opinion in Neurobiology} \textbf{6}(2), 171--178 (1996).

\bibitem{singer1999}
W.~Singer, ``Neuronal synchrony: A versatile code for the definition of relations?'' \emph{Neuron} \textbf{24}(1), 49--65 (1999).

\bibitem{varela2001}
F.~Varela, J.-P.~Lachaux, E.~Rodriguez, and J.~Martinerie, ``The brainweb: Phase synchronization and large-scale integration,'' \emph{Nature Reviews Neuroscience} \textbf{2}(4), 229--239 (2001).

\bibitem{fries2015}
P.~Fries, ``Rhythms for cognition: Communication through coherence,'' \emph{Neuron} \textbf{88}(1), 220--235 (2015).

\bibitem{casali2013}
A.~G.~Casali et al., ``A theoretically based index of consciousness independent of sensory processing and behavior,'' \emph{Science Translational Medicine} \textbf{5}(198), 198ra105 (2013).

\bibitem{mashour2020}
G.~A.~Mashour, P.~Roelfsema, J.-P.~Changeux, and S.~Dehaene, ``Conscious processing and the global neuronal workspace hypothesis,'' \emph{Neuron} \textbf{105}(5), 776--798 (2020).

\bibitem{lee2013}
U.~Lee et al., ``Disruption of frontal--parietal communication by ketamine, propofol, and sevoflurane,'' \emph{Anesthesiology} \textbf{118}(6), 1264--1275 (2013).

\bibitem{dehaene2011}
S.~Dehaene and J.-P.~Changeux, ``Experimental and theoretical approaches to conscious processing,'' \emph{Neuron} \textbf{70}(2), 200--227 (2011).

\bibitem{tononi2016}
G.~Tononi, M.~Boly, M.~Massimini, and C.~Koch, ``Integrated information theory: From consciousness to its physical substrate,'' \emph{Nature Reviews Neuroscience} \textbf{17}(7), 450--461 (2016).

\bibitem{hameroff2014}
S.~Hameroff and R.~Penrose, ``Consciousness in the universe: A review of the `Orch OR' theory,'' \emph{Physics of Life Reviews} \textbf{11}(1), 39--78 (2014).

\bibitem{tegmark2000}
M.~Tegmark, ``Importance of quantum decoherence in brain processes,'' \emph{Physical Review E} \textbf{61}(4), 4194--4206 (2000).

\bibitem{caldeira1983}
A.~O.~Caldeira and A.~J.~Leggett, ``Path integral approach to quantum Brownian motion,'' \emph{Physica A} \textbf{121}(3), 587--616 (1983).

\bibitem{penrose1996}
R.~Penrose, ``On gravity's role in quantum state reduction,'' \emph{General Relativity and Gravitation} \textbf{28}(5), 581--600 (1996).

\bibitem{diosi1989}
L.~Diósi, ``Models for universal reduction of macroscopic quantum fluctuations,'' \emph{Physical Review A} \textbf{40}(3), 1165--1174 (1989).

\end{thebibliography}

\end{document}
